\section{Theoretical Background}
\label{sec:background}

Whereas Section~\ref{sec:related_work} compares empirical studies, this section
defines the physical target, the forecasting task, and the mechanisms of
TimesNet and DilatedRNN. The distinction is important because neither an
architecture name nor a nominal horizon uniquely specifies a forecasting
experiment.

\subsection{Global horizontal irradiance and temporal support}

Global horizontal irradiance (GHI) is the instantaneous solar radiant flux
received per unit area by a horizontal surface. For the sun above the horizon,
\begin{equation}
  \mathrm{GHI}
  =
  \mathrm{DHI}
  +
  \mathrm{DNI}\cos(\theta_z),
  \label{eq:ghi}
\end{equation}
where DHI is diffuse horizontal irradiance, DNI is direct normal irradiance,
and $\theta_z$ is the solar zenith angle \cite{duffie2013}. Irradiance is a
power density, conventionally expressed in W/m$^2$; time-integrated irradiation
is an energy density and is a different quantity.

The NSRDB product used in this study estimates the radiative components from
satellite observations and physical models \cite{sengupta2018}. A value at the
grid cell nearest a factory coordinate is consequently a model-based estimate
for that location, not a reading from a sensor mounted at the factory. This
distinction also precludes interpreting GHI directly as photovoltaic power,
which would require information about module area, orientation, efficiency,
temperature, losses, and plant operation.

Let $y_{\ell,t}^{(1\mathrm{h})}$ be hourly mean GHI for location $\ell$. The
mean irradiance at a coarser temporal support $\Delta$ is
\begin{equation}
  y_{\ell,\tau}^{(\Delta)}
  =
  \frac{1}{N_{\ell,\tau}}
  \sum_{t\in\mathcal{I}_{\ell,\tau}^{(\Delta)}}
  y_{\ell,t}^{(1\mathrm{h})},
  \label{eq:temporal_aggregation}
\end{equation}
where $\mathcal{I}_{\ell,\tau}^{(\Delta)}$ contains the valid hourly samples in
interval $\tau$ and $N_{\ell,\tau}$ is their number. Equation
\eqref{eq:temporal_aggregation} preserves W/m$^2$ and defines a daily or
monthly mean; it is not an accumulated daily or monthly irradiation. Coverage
requirements must be applied before aggregation so that intervals with missing
hours are not treated as complete.

\subsection{Forecast resolution, context, and horizon}

For a temporal resolution $\Delta$, let
$\mathbf{X}_{\ell,t}\in\mathbb{R}^{C\times P}$ contain the $C$ observations
and $P$ admissible variables observed strictly before target interval $t$.
Throughout the article, the origin label $t$ identifies the first predicted
interval; its value is unavailable to the forecaster. A deterministic
multi-output forecaster is a mapping
\begin{equation}
  \widehat{\mathbf{y}}_{\ell,t:t+H-1}
  =
  f_{\theta}\!\left(\mathbf{X}_{\ell,t},\mathbf{z}_{\ell}\right)
  \in\mathbb{R}^{H},
  \label{eq:forecast_map}
\end{equation}
where $\mathbf{z}_{\ell}$ denotes static location information when used,
$C$ is the context length, and $H$ is the number of predicted intervals.
For fixed-duration hourly and daily intervals, the corresponding spans are
$C\Delta$ and $H\Delta$; monthly spans follow civil months of unequal length.
Thus, 72 hourly outputs and three daily outputs have the same nominal
duration but different targets and information content. Likewise, predicting
one monthly mean is a one-step monthly forecast, not a multi-month trajectory
\cite{diagne2013,yang2018}.

In a direct multi-output strategy, all $H$ values in
\eqref{eq:forecast_map} are produced jointly. A recursive strategy instead
feeds one-step predictions back as inputs and can accumulate error with
horizon \cite{guariso2020}. The strategy must therefore be held constant in a
model comparison. A local model estimates separate parameters for each
location; a global model shares $\theta$ across locations and may condition on
$\mathbf{z}_{\ell}$ \cite{montero2021}. Sharing can increase the effective
training sample, but it does not guarantee equal local performance.

Persistence, diurnal or seasonal naive forecasts, and simple autoregressive
references expose how much skill is attributable to continuity and
seasonality rather than architectural complexity. Their suitability depends
on resolution and horizon, so more than one transparent reference may be
needed \cite{voyant2022}.

\subsection{TimesNet}

TimesNet models one-dimensional sequences through a set of
period-conditioned two-dimensional representations \cite{wu2023}. Let
$\mathbf{X}\in\mathbb{R}^{B\times T\times C_x}$ be a batch processed by a
TimesBlock, and let $F_{b,f,c}$ denote the temporal fast Fourier transform
(FFT) of sample $b$ at frequency bin $f$ and channel $c$. The compact variant
evaluated here selects periods independently for each sample, averaging
amplitudes only over that sample's channels,
\begin{equation}
  A_{b,f}
  =
  \frac{1}{C_x}
  \sum_{c=1}^{C_x}\left|F_{b,f,c}\right|.
  \label{eq:timesnet_frequency}
\end{equation}
The zero-frequency component is excluded, and each sample independently uses
its $k$ largest amplitudes to select frequency indices $f_{b,i}$ and integer
periods
\begin{equation}
  \{f_{b,1},\ldots,f_{b,k}\}
  =
  \operatorname*{arg\,top\text{-}k}_{1\leq f\leq\lfloor T/2\rfloor}
  A_{b,f},
  \qquad
  p_{b,i}=\left\lfloor\frac{T}{f_{b,i}}\right\rfloor .
  \label{eq:timesnet_periods}
\end{equation}
The retained aggregation amplitude is therefore
$a_{b,i}=A_{b,f_{b,i}}$. This sample-specific selection ensures that changing
other windows in a computational batch cannot change sample $b$'s periods.

For each $p_{b,i}$, the sequence is zero-padded to
$L_{b,i}=p_{b,i}\lceil T/p_{b,i}\rceil$ and reshaped so that one axis indexes
position within a period and the other indexes successive periods. A
two-dimensional Inception-style convolution then learns intraperiod and
interperiod variations. If $\mathcal{G}_{b,i}$ is the processed tensor and
$\mathcal{R}_{b,i}^{-1}$ restores chronological order, the period-specific
outputs are combined with sample-specific amplitude-derived weights
\begin{equation}
\begin{aligned}
  \alpha_{b,i}
  &=
  \frac{\exp(a_{b,i})}
       {\sum_{j=1}^{k}\exp(a_{b,j})},
  \\
  \widetilde{\mathbf{X}}_b
  &=
  \mathbf{X}_b
  +
  \sum_{i=1}^{k}
  \alpha_{b,i}
  \left[
    \mathcal{R}_{b,i}^{-1}(\mathcal{G}_{b,i})
  \right]_{1:T}.
\end{aligned}
  \label{eq:timesnet_aggregation}
\end{equation}
The residual connection preserves the original representation. Before value
embedding, the implementation also standardizes each input window by its own
temporal mean and population standard deviation, with $10^{-5}$ added to the
variance, and reverses that transformation at the output. This internal,
causal window standardization is applied after the external site-wise min--max
transformation described in Section~\ref{sec:data}. Stacked TimesBlocks and a
task-specific projection produce the forecast. The two-dimensional tensors are
temporal reorganizations, not geographic images; their axes express
within-period position and recurrence across periods. Period selection and
aggregation weights are both sample-specific; batch composition is therefore
a computational detail rather than part of the forecast information set.
TimesNet has also been combined with signal decomposition for week-ahead
hourly solar forecasting \cite{zhao2024timesnet}, but decomposition is not an
intrinsic component of the TimesNet architecture.

\subsection{Dilated recurrent neural networks}

A DilatedRNN changes the recurrent connection pattern rather than reshaping the
sequence \cite{chang2017}. Let $\mathbf{c}_{t}^{(q)}$ denote the recurrent
state in layer $q$, and let $\mathbf{x}_{t}^{(q)}$ be the layer input, with
$\mathbf{x}_{t}^{(1)}=\mathbf{x}_t$ and
$\mathbf{x}_{t}^{(q)}=\mathbf{c}_{t}^{(q-1)}$ for $q>1$. A dilated recurrent
layer updates
\begin{equation}
  \mathbf{c}_{t}^{(q)}
  =
  \phi_q\!\left(
    \mathbf{x}_{t}^{(q)},
    \mathbf{c}_{t-d_q}^{(q)}
  \right),
  \label{eq:dilated_rnn}
\end{equation}
where $\phi_q$ can be an ordinary RNN, gated recurrent unit, or LSTM cell
\cite{chang2017,hochreiter1997}. The recurrent edge from $t-1$ is replaced by the edge
from $t-d_q$ in that layer; dilation is therefore not an additional residual
connection.

With an integer base $M>1$, the standard exponentially increasing schedule is
\begin{equation}
  d_q=M^{q-1},\qquad q=1,\ldots,Q.
  \label{eq:dilation_schedule}
\end{equation}
The first layer retains local recurrence when $d_1=1$, while higher layers
connect states at progressively coarser temporal separations. This schedule
shortens recurrent paths between distant observations and exposes multiple
temporal scales. It does not discard or downsample the intervening
observations: the interleaved recurrent chains are recombined at their
original time indices. An output head maps the resulting states to one or
several forecast steps.

TimesNet and DilatedRNN consequently encode multi-scale dependence through
different inductive biases. TimesNet discovers dominant Fourier periods and
applies two-dimensional convolutions; DilatedRNN fixes a hierarchy of sparse
recurrent lags. Neither architecture, by itself, defines the context,
forecasting strategy, horizon, or evaluation protocol. Those elements must be
specified and matched before the implementations can be compared.
Even then, unequal capacity and optimization responses limit attribution to
any individual architectural mechanism.

\subsection{Atmospheric variability and climate context}

Solar geometry creates a deterministic diurnal envelope, whereas clouds and
other atmospheric constituents modulate the irradiance reaching the surface.
Cloud advection can create rapid GHI changes and photovoltaic power ramps
\cite{chow2011,chen2020cloudramps}. Aggregation reduces some high-frequency
variation but does not remove persistent cloudy regimes, synoptic variability,
or seasonality.

Low GHI alone is not evidence of precipitation: the value may reflect solar
elevation, cloud optical properties, aerosols, or the temporal aggregation
window. Descriptions such as rainy or cloudy therefore require an independent
meteorological field with aligned time and location. Reanalysis products can
provide such contextual fields, subject to their own spatial resolution and
model uncertainty \cite{hersbach2020,copernicus2026era5}.

K{\"o}ppen--Geiger classes summarize long-term temperature and precipitation
regimes \cite{beck2018}. They are useful for describing the geographic
diversity of the sampled locations, but they are not dynamic forecast inputs
unless explicitly encoded. With ten sites, differences between classes are
exploratory associations rather than causal estimates of climate effects on
forecast error.
